\subsection{Proof of Theorem 1}
\label{app:proof-linlora-retention}

\begin{proof}
Recall that
\[
W^{\ell \ell}
=
U^{pre}M^{\ell \ell},
\qquad
M^{\ell \ell}
=
\begin{bmatrix}
M'^{pre} \\
\Delta M^{\ell \ell}
\end{bmatrix}.
\]
Therefore,
\[
W^{\ell \ell}h_{i-1}
=
U^{pre}M^{\ell \ell}h_{i-1}
=
U^{pre}
\begin{bmatrix}
M'^{pre} \\
\Delta M^{\ell \ell}
\end{bmatrix}
h_{i-1}.
\]
Let
\[
M^{\ell \ell}h_{i-1}
=
\begin{bmatrix}
z^{pre}_1 \\
\vdots \\
z^{pre}_R \\
z^{ft}_{R+1} \\
\vdots \\
z^{ft}_{R+r}
\end{bmatrix}.
\]
By the block structure and linearity of the matrix-vector product,
\[
W^{\ell \ell}h_{i-1}
=
U^{pre}
\left(
\begin{bmatrix}
z^{pre}_1 \\
\vdots \\
z^{pre}_R \\
0 \\
\vdots \\
0
\end{bmatrix}
+
\begin{bmatrix}
0 \\
\vdots \\
0 \\
z^{ft}_{R+1} \\
\vdots \\
z^{ft}_{R+r}
\end{bmatrix}
\right).
\]
Since the columns of $U^{pre}$ are orthonormal, the two sums span disjoint
orthogonal subspaces. Hence,
\[
W^{\ell \ell}h_{i-1}
=
\sum\limits_{j=1}^{R} z^{pre}_j U^{pre}_j
+
\sum\limits_{j=R+1}^{R+r} z^{ft}_j U^{pre}_j .
\]
This establishes that the pretrained and fine-tuned components are represented
in orthogonal blocks of the fixed pretrained left-singular basis.
\end{proof}


\subsection{Proof of Proposition 2}
\label{app:proof-proposition-2-scaler-induces-mixing-bounds}
\begin{proof}
The first block is
\[
C_{11}
=
(U_R^\top E \bar{U}_r)\Sigma'_r(\bar{V}_r^\top D V_R).
\]
Using submultiplicativity of the operator norm,
\[
\|C_{11}\|_2
\leq
\|U_R^\top E \bar{U}_r\|_2
\|\Sigma'_r\|_2
\|\bar{V}_r^\top D V_R\|_2 .
\]
Since $D$ is diagonal and therefore symmetric,
\[
\|\bar{V}_r^\top D V_R\|_2
=
\|V_R^\top D \bar{V}_r\|_2
=
\nu_D .
\]
Thus,
\[
\|C_{11}\|_2
\leq
\mu_E \|\Sigma'_r\|_2 \nu_D .
\]
The other bounds follow identically from submultiplicativity, using
\[
\|\bar{V}_r^\top D V_r\|_2 \leq \|D\|_2,
\qquad
\|U_r^\top E \bar{U}_r\|_2 \leq \|E\|_2 .
\]
\end{proof}


\subsection{Proof of Theorem 2}
\label{app:proof-thm-2-approximate-tail-preservation}
\begin{proof}
By the block decomposition,
\[
C-C_{\mathrm{tail}}
=
\begin{bmatrix}
C_{11} & C_{12} \\
C_{21} & 0
\end{bmatrix}.
\]
Therefore,
\[
\|C-C_{\mathrm{tail}}\|_2
\leq
\sqrt{
\|C_{11}\|_2^2
+
\|C_{12}\|_2^2
+
\|C_{21}\|_2^2
}.
\]
Using Proposition 2,
\[
\|C-C_{\mathrm{tail}}\|_2
\leq
\|\Sigma'_r\|_2
\sqrt{
(\mu_E\nu_D)^2
+
(\mu_E\|D\|_2)^2
+
(\|E\|_2\nu_D)^2
}.
\]
If $\mu_E\to 0$ and $\nu_D\to 0$, while $\|E\|_2$, $\|D\|_2$, and
$\|\Sigma'_r\|_2$ remain bounded, the right-hand side converges to zero.
\end{proof}

